% METADATA DEFINITIONS (Hidden from PDF output)












\documentclass[aspectratio=169, 12pt]{beamer}

\usepackage{fontspec}
\usepackage{polyglossia}
\usepackage{tikz}
\usepackage{xcolor}
\usepackage{graphicx}
\usepackage{amsmath, amssymb}
\usepackage{listings}
\usepackage{float}
\usepackage[object = vectorian]{pgfornament}

% --- Beamer Theme Configuration ---
\usetheme{Berlin}
\usecolortheme{default}
\usefonttheme{default}

% Metropolis specific tweaks


% Devanagari Support logic

    \setmainfont[Script=Devanagari, AutoFakeBold=1.8]{Tiro Devanagari Sanskrit}
    \setsansfont[Script=Devanagari, AutoFakeBold=1.8]{Tiro Devanagari Sanskrit}


% --- Devanagari Digit Macros ---
\makeatletter
\def\devanagaridigits#1{\expandafter\@devanagaridigits\number#1@}
\def\@devanagaridigits#1{%
  \ifx#1@\else
    \ifcase#1 ०\or १\or २\or ३\or ४\or ५\or ६\or ७\or ८\or ९\fi
    \expandafter\@devanagaridigits
  \fi
}
\makeatother


    \renewcommand{\thepage}{\devanagaridigits{\value{page}}}
    \setbeamertemplate{page number in head/foot}[totalnumber]


% --- Navigation Symbols ---

    \setbeamertemplate{navigation symbols}{}


% --- Section Titles frame logic ---

    \AtBeginSection[]{
      \begin{frame}
      \vfill
      \centering
      \begin{beamercolorbox}[sep=8pt,center,shadow=true,rounded=true]{title}
        \usebeamerfont{title}\insertsectionhead\par%
      \end{beamercolorbox}
      \vfill
      \end{frame}
    }


% --- Custom Macros ---
\newcommand{\maketitleframe}[4]{
    \title{#1}
    \subtitle{#2}
    \author{#3}
    \date{#4}
    \begin{frame}
        \titlepage
    \end{frame}
}

\newcommand{\thankyouframe}[2]{
    \begin{frame}
        \vfill
        \centering
        \begin{tikzpicture}
            \node[inner sep=20pt, fill=structure!10, rounded corners=15pt, draw=structure, thick] (box) {
                \begin{minipage}{0.6\textwidth}
                    \centering
                    {\Huge \bfseries \color{structure} #1} \par
                    \vspace{0.5cm}
                    {\large #2}
                \end{minipage}
            };
        \end{tikzpicture}
        \vfill
    \end{frame}
}

% --- Alignment and Spacing Macros ---
\newcommand{\vspacer}[1]{\vspace{#1 cm}}

% --- Shloka Environment ---
\newcounter{shlokaline}
\newenvironment{shloka}[1][alternating]{
    \par\vspace{0.3cm}
    \begingroup
    \parindent=0pt
    \def\shlokaalign{#1}
    \setcounter{shlokaline}{0}
    \ifthenelse{\equal{\shlokaalign}{alternating}}{
        \let\oldpar\par
        \def\par{\oldpar\stepcounter{shlokaline}}
        \everypar{\ifodd\value{shlokaline}\raggedleft\else\raggedright\fi}
    }{
        \ifthenelse{\equal{\shlokaalign}{left}}{\raggedright}{}
        \ifthenelse{\equal{\shlokaalign}{right}}{\raggedleft}{}
        \ifthenelse{\equal{\shlokaalign}{center}}{\centering}{}
    }
}{
    \par\endgroup
    \vspace{0.3cm}
}

% --- Vishesh Shloka Environment (with Ornament) ---
\newenvironment{visheshshloka}[1][alternating]{
    \begin{shloka}[#1]
}{
    \end{shloka}
    \vspace{-0.4cm}
    \begin{center}
        \color{theme}\pgfornament[width=0.2\textwidth]{88}
    \end{center}
    \vspace{0.1cm}
}

% --- Vishesh Shloka Two Environment (with Custom Ornament) ---
\newenvironment{visheshshlokatwo}[2][alternating]{
    \def\visheshornament{#2}
    \begin{shloka}[#1]
}{
    \end{shloka}
    \vspace{-0.4cm}
    \begin{center}
        \color{theme}\pgfornament[width=0.2\textwidth]{\visheshornament}
    \end{center}
    \vspace{0.1cm}
}

% --- Image Macros ---
\newcommand{\customimage}[3]{
    \begin{figure}[H]
        \ifthenelse{\equal{#2}{Center}}{\centering}{}
        \ifthenelse{\equal{#2}{Left}}{\raggedright}{}
        \ifthenelse{\equal{#2}{Right}}{\raggedleft}{}
        \includegraphics[width=#3\textwidth]{#1}
    \end{figure}
}

% --- Vishesh Butta (Repeated Ornaments) ---
\newcommand{\visheshbutta}[5]{%
    \par\vspace{0.2cm}
    \noindent\begin{minipage}{\textwidth}
        \ifthenelse{\equal{#1}{Center}}{\centering}{}%
        \ifthenelse{\equal{#1}{Left}}{\raggedright}{}%
        \ifthenelse{\equal{#1}{Right}}{\raggedleft}{}%
        \foreach \n in {1,...,#3}{%
            \pgfornament[width=#4]{#2}\hspace{#5}%
        }%
    \end{minipage}\par\vspace{0.2cm}
}

% --- Custom Colors ---
\definecolor{theme}{named}{blue}

\begin{document}

\maketitleframe{विभक्ति छुट्ट्याऊ आन्दोलन}{विभक्ति छुट्ट्याउँदा राम्रो कि नछुट्ट्याउँदा ?}{क्वैराला संस्कृतः}{चैत्र, २०८२}

\begin{frame}{विषय सूची}\tableofcontents\end{frame}

\begin{frame}[fragile]{प्रारम्भ}{मुद्दा के हो ?}

\begin{alertblock}{प्रश्न}

`सबैको घरमा' हो कि `सबै को घर मा' ?

विशेषतः विभक्ति छुट्ट्याएर लेख्दा राम्रो कि जोडेर ?

\end{alertblock}

\begin{itemize}
\item केही नेपाली लेखकहरूले विभक्ति छुट्ट्याइएको देखिन्छ । त्यो प्रयोगको पृष्ठमा के छ ?


\item मानक नेपालीमा यथावत् रहेको विभक्ति जोडिने नियम कत्तिको ठिक छ ?


\end{itemize}
\end{frame}

\section{पृष्ठभूमि}

\begin{frame}[fragile]{भाषाका दुई स्वभाव}{संश्लेेषणात्मक र विश्लेषणात्मक}

\begin{columns}[t]

\begin{column}{0.5\textwidth}

\textbf{संश्लेषणात्मक स्वभाव}

\begin{itemize}
\item एकै शब्दले धेरै कुरा व्यक्त गर्ने ।


\end{itemize}
\begin{alertblock}{उदा०}

पशुहरूको ध्वनि ।

पिपृच्छामि ।

नेपाली = खान्छु ।

\end{alertblock}

\end{column}

\begin{column}{0.5\textwidth}

\textbf{विश्लेषणात्मक स्वभाव}

\begin{itemize}
\item धेरै शब्दहरूद्वारा भावहरू व्यक्त गरिने ।


\end{itemize}
\begin{exampleblock}{उदा०}

मान्छेहरूको भाषा ।

अहं प्रष्टुम् इच्छामि = अहं प्रश्नं कर्तुम् इच्छामि ।

अङ्ग्रेजी = I eat.

\end{exampleblock}

\end{column}

\end{columns}

\end{frame}

\begin{frame}[fragile]{प्राकृत > नेपाली}{नेपाली कता पर्दछ?}

\begin{itemize}
\item नेपालीमा केही संश्लेषणात्मक र केही विश्लेषणात्मक विशेषताहरू छन् । (अर्को कुनै वेला कुरा गरौँला)


\item केही नेपाली संश्लेषणात्मक विशेषताहरू संस्कृत पछि पनि विकास भएका छन् ।


\item संस्कृतको जस्तो धातुको स्तरमा गएर मानकीकरण भएजस्तो विश्लेषण नेपाली भाषाविज्ञानमा छैन ।


\end{itemize}
\end{frame}

\section{विभक्तिहरूको प्राचीन विकासक्रम}

\begin{frame}[fragile]{विभक्ति के हो ?}{}

\begin{itemize}
\item विभक्ति एउटा परस्थानिक (Postposition) हो । कुनै पनि संज्ञाको पछाडि आउने शब्दांश, जसले संज्ञाको कारक/वचन इत्यादिलाई बुझाउँछ परस्थानिक भनिन्छ । अङ्ग्रेजीमा पूर्वस्थानिक (Preposition) हुन्छ ।


\end{itemize}
\begin{exampleblock}{उदा०}

घर\textit{मा}, घर\textit{भित्र} इत्यादि ‍ = परस्थानिक

On the table, Under the table = पूर्वस्थानिक

\end{exampleblock}

\begin{itemize}
\item नेपालीमा शीर्ष दुई परस्थानिकहरू छन् । नामयोगी र विभक्ति चिह्न ।


\end{itemize}
\end{frame}

\begin{frame}[fragile]{नेपाली विभक्ति चिह्न विकासक्रम}{}

\begin{itemize}
\item ``मा'' = `मध्य' शब्दबाट आएको हो। (मध्य > मज्झ > माझ > मा)


\item ``बाट'' = `वर्त्मन्' (बाटो) बाट आएको हो। (वर्त्मन् > वट्ट > बाटो > बाट)


\item ``देखि'' = `दृश्' धातुको पूर्वकालिक रूपबाट। (दृष्टि/दृक्ष्य > देखिअ > देखि)


\item ``ले'' = `लग्' धातुबाट। (लग्न > लग्गिअ > लागि > ले)


\item ``लाई'' = `लग्' धातुबाटै विकसित। (लग्न > लग्गिअ > लागि > लाई)


\item ``को/का/की'' = `कृत' शब्दबाट आएको हो। (कृतक > केरअ > केरा > को)


\item ``तिर'' = `तीर' शब्दबाट। (तीर > तीर > तीर > तिर)


\item ``सँग'' = `सङ्ग' शब्दबाट। (सङ्ग > सङ्ग > सङ्ग > सँग)


\item ``सित'' = `सहित' शब्दबाट। (सहित > सहिअ > सित)


\item ``सम्म'' = `सम' शब्दबाट। (सम > सम्म > सम्म)


\item ``द्वारा'' = `द्वार्' शब्दबाट। (द्वारा (द्वार् ले) > दुआर > द्वारा)


\end{itemize}
\end{frame}

\begin{frame}[fragile]{व्याकरणिकीरण (Grammaticalization)}{}

\begin{block}{परिभाषा}

व्याकरणिकीकरण एउटा यस्तो भाषाई प्रक्रिया हो, जसमा कुनै स्वतन्त्र र ठोस अर्थ बोक्ने शब्द (Lexical Item) समयक्रमसँगै आफ्नो मूल अर्थ गुमाउँदै व्याकरणको एउटा नियम, प्रत्यय वा सहायक तत्त्व (Grammatical/Functional Item) मा परिणत हुन्छ।

\end{block}

जस्तै-

अङ्ग्रेजीमा `will' शब्द पहिले `willan' भन्ने क्रियापद जसले `चाहनु' भन्ने बुझाउँछ, त्यो आधुनिक अङ्ग्रेजीमा सहायक क्रियापद (Auxiliary Verb) भएको छ ।

\end{frame}

\begin{frame}[fragile]{पद-वियोगको समर्थन}{}

\begin{itemize}
\item मूल भाषामा भिन्न शब्दका  रूपमा स्थापित भएकाले


\item ती चिह्नहरू शब्दरूप नभएकाले (संस्कृतमा जस्तो `ले' र `मा' शब्दानुसार परिवर्तन हुँदैन)


\end{itemize}
\begin{block}{नोट}

सर्वनाममा भने केही स्थानमा रूप परिवर्तन हुन्छ, जस्तै (म + को = मेरो, ऊ + ले = उसले/उल्ले, आफू + को = आफ्नो, हामी + को = हाम्रो)  । त्यसैले पद-वियोग गर्ने जमातले सर्वनाम भने योग गरेर लेख्दछन् । (हिन्दी सरह)

\end{block}

\begin{itemize}
\item केही सर्वनामका रूपहरू वास्तवमा मूलबाट एकै चोटि अपभ्रंश भएका कारणले छुट्टै शब्द हो भनेर पनि भन्न सकिन्छ । (अस्माकम् > अम्हाक/अम्हार > हाम्रो)


\end{itemize}
\end{frame}

\begin{frame}[fragile]{मध्यकालीन नेपाली}{}

मध्यकालीन नेपाली अभिलेखहरूमा केही अतिरिक्त परस्थानिकहरू छन् । जस्तै -

\begin{itemize}
\item हर = पछि (खाइ हर, अहिले  `खाएर' भनेर संश्लेषण भएको छ । `खाए पछि' को `खाएसी' जस्तै)


\item थि/थ्यौ/शुँ = सँग (अन्नशुँ/हामीथ्यौ/हामुथि)


\item सुदा = समेत (षेतसुदा, अहिलेको `सुद्ध' जस्तै)


\item छयो = सँग (कार्कीसँग)


\end{itemize}
यस्ता प्रयोगले भाषाको संश्लेषणात्मक हुन खोजेको यत्न र जोडेरै लेख्दा उचित हुन्छ कि जस्तो देखिन्छ ।

\end{frame}

\section{हिन्दी र अन्य भाषाहरूसँग तुलना}

\begin{frame}[fragile]{हिन्दी vs नेपाली}{हिन्दी र नेपालीको स्वभावमा फरक}

क्रियापदको विषयमा हिन्दी भाषा नेपालीभन्दा धेरै विश्लेषणात्मक भाषा हो। जस्तै -

\begin{itemize}
\item `जाता हुं' (हिन्दी) = `जान्छु' (नेपाली)


\item `नहीं जाता हुं' (हिन्दी) = `जान्न/जादिनँ' (नेपाली) [यो विषयमा नेपाली संस्कृतभन्दा पनि संश्लेषणात्मक छ]


\item `अरे ! राम तो घर चला गया' / `राम तो घर चला गया' (हिन्दी) = राम घर गएछ (नेपाली)


\end{itemize}
\end{frame}

\begin{frame}[fragile]{हिन्दी vs नेपाली}{हिन्दी र नेपालीको स्वभावमा समानता}

तर विभक्तिको विषयमा भने नेपाली र हिन्दीमा (तथा अङ्ग्रेजीमा पनि) धेरै फरक नभई वस्तुतः विश्लेषणात्मक छन् । जस्तै -

\begin{itemize}
\item `घर में रखे टेबल के नींचे से चाकु निकाल कर ला' (हिन्दी) = `घरमा राखेको टेबलको मुनिबाट चाकु निकालेर ले' (नेपाली) वा छुट्ट्याएर लेख्दा - `घर मा राखेको टेबल को मुनि बाट चाकु निकालेर ले' = home at kept table below from knife take out and get = Take out and get the knife from under the table kept at home. (अङ्ग्रेजी)


\item अङ्ग्रेजीको `Preposition' ‍=  नेपाली/हिन्दीको आसन्न `Postposition'


\item केही भाषाविज्ञानीहरूले हिन्दीमा र तदनुसार नेपालीमा कारक २/३ मात्र भएको बताउँछन् । सरल, तिर्यक् र सम्बोधन


\end{itemize}
\end{frame}

\begin{frame}[fragile]{अन्य भाषामा के गरिन्छ ?}{}

\begin{itemize}
\item हिन्दी वा हिन्दीको अतिप्रभाव भएका भाषाहरूभन्दा बाहेक जोडेर लेखिन्छ । मराठीमा `रामाचा', गुजरातीमा `घरमां', बङ्गालीमा `रामेर', उडीयामा `रामकु' इत्यादि ।


\item पारसीको अतिप्रभावले रसम-उल-ख़तमा हिन्दवी भाषा बढी प्रचलित भएसँगै विश्लिष्ट रूपमा लेख्ने चलन भएको देखिन्छ ।


\end{itemize}
\end{frame}

\section{संज्ञानात्मक विलम्ब}



\end{document}